\chapter*{Conclusiones y recomendaciones}
\addcontentsline{toc}{chapter}{Conclusiones y recomendaciones}
\markboth{Conclusiones y recomendaciones}{Conclusiones y recomendaciones}

\section*{Conclusiones}

El presente trabajo cumplió con el objetivo general planteado: se desarrolló \emph{EduFEM}, un software educativo de escritorio escrito en el lenguaje de programación Python sobre bibliotecas numéricas de código abierto, para el análisis por el Método de los Elementos Finitos (MEF) de problemas bidimensionales de elasticidad lineal en tensión plana y deformación plana, con elementos cuadriláteros isoparamétricos Q4 y Q9. La herramienta integra, en un único entorno coherente, las tres dimensiones que se propusieron como ejes del proyecto: la formulación matemática del método, la construcción y visualización interactiva del modelo, y la verificación numérica de los resultados. El trabajo regresa así a la doble brecha que lo originó —una dificultad de aprendizaje genuina, ligada a la naturaleza abstracta del método, y la insuficiencia de herramientas que la atiendan exponiendo el proceso de cálculo de manera transparente y guiada— y a la interrogante que planteó como problema: de qué manera la exposición de ese canal apoya la comprensión de los fundamentos del método. Como primer ciclo de una investigación basada en diseño, el trabajo responde la parte de esa interrogante que puede establecerse sobre el artefacto —que la exposición es relevante, correcta, consistente con los principios de aprendizaje y practicable— y deja diseñada la prueba de campo que responderá la parte que reside en el estudiante.

A continuación se presentan las conclusiones específicas, ordenadas en correspondencia con los objetivos planteados y con el capítulo que desarrolla cada uno.

En relación con la fundamentación teórica (\autoref{cap:marco_teorico}), se estableció la base que el motor, los módulos educativos, la memoria de cálculo y la evaluación comparten: la formulación del MEF para la elasticidad lineal plana con elementos isoparamétricos Q4 y Q9 —de la formulación variacional a la recuperación de tensiones—, la calidad de malla y los criterios de verificación y validación; los diez principios de aprendizaje, generales y específicos de la enseñanza del método, que la exposición del canal de cálculo encarna; y los criterios de calidad de una intervención educativa, el mapa de conjeturas y los instrumentos de evaluación sin usuarios con que la investigación basada en diseño organiza el contraste. Esa base no es un resultado medido sino el fundamento del trabajo; lo que los resultados comprueban es que el motor construido sobre ella se comporta como la teoría predice, que las mismas expresiones que el capítulo expone son las que el código implementa, los módulos exhiben y la memoria sustituye numéricamente, y que cada decisión de diseño remite a un principio documentado.

En relación con la implementación del motor de cálculo (\autoref{cap:diseno}; verificación en el \autoref{cap:resultados}), se construyó un núcleo numérico completo en \texttt{NumPy} y \texttt{SciPy} que cubre la totalidad del canal de cálculo del MEF para elementos isoparamétricos: funciones de forma bilineales (Q4) y bicuadráticas (Q9), evaluación del Jacobiano, ensamblaje de la matriz de deformación-desplazamiento $\bm{B}$, construcción de la matriz constitutiva $\bm{D}$ para ambos estados planos, integración de la matriz de rigidez elemental por cuadratura de Gauss, ensamblaje global, imposición de restricciones (incluyendo desplazamientos prescritos no homogéneos) y recuperación de tensiones por extrapolación desde los puntos de Gauss. La corrección de este motor quedó verificada mediante el método de soluciones manufacturadas (MMS) en tensión plana y en deformación plana, sobre mallas regulares y distorsionadas: las tasas de convergencia asintóticas observadas coinciden con las predicciones teóricas $\mathcal{O}(h^{p+1})$ en norma $L^2$ y $\mathcal{O}(h^{p})$ en seminorma $H^1$, esto es, $\mathcal{O}(h^{2{,}00})$ y $\mathcal{O}(h^{1{,}00})$ para Q4, y $\mathcal{O}(h^{3{,}00})$ y $\mathcal{O}(h^{2{,}00})$ para Q9, y el campo de tensiones recuperado converge a $\mathcal{O}(h^{1{,}5})$ en Q4 y $\mathcal{O}(h^{2})$ en Q9, con una matriz de extrapolación que reproduce exactamente los campos polinómicos de cada elemento. La obtención de estas tasas constituye evidencia rigurosa de que el motor está libre de errores de orden en la formulación de los elementos y en la recuperación de tensiones.

Respecto al diseño de la interfaz gráfica (\autoref{cap:diseno}; resultado en la \autoref{sec:cobertura}), se logró una aplicación organizada en las tres fases canónicas del análisis —pre-proceso, proceso y post-proceso— que comparten un único lienzo interactivo: en los tres casos de referencia del \autoref{cap:resultados} el modelo se construyó o importó, se resolvió y se inspeccionó sin cambiar de vista, y la \autoref{tab:fases} registra qué operaciones ofrece cada fase y con qué instrumento expone cada etapa del canal. Esta organización permite que el estudiante construya la geometría, asigne materiales, cargas y restricciones, resuelva el sistema y visualice los campos de desplazamiento y tensión sin cambiar de contexto visual, con selección bidireccional entre las tablas de datos y el lienzo, deshacer y rehacer sobre el modelo completo y validación de la salud del modelo previa a la resolución. La evaluación heurística la juzgó cumplida en diseño de la presentación y reusabilidad, y parcial en usabilidad de la interacción y accesibilidad, dimensiones cuya evidencia definitiva requiere usuarios; la orientación minimalista de la interfaz es una decisión de diseño de este trabajo, trazable al principio de gestión de la carga cognitiva.

En cuanto a los módulos educativos (\autoref{cap:diseno}; cobertura en la \autoref{sec:cobertura}), se desarrollaron ocho módulos que exponen, paso a paso y sobre el modelo real que el estudiante ha construido, las etapas del canal de cálculo del MEF desde el mapeo isoparamétrico hasta el ensamblaje global (M1 a M7, \autoref{tab:modulos}), más la calidad de malla en el pre-proceso (M0); la solución del sistema y la recuperación de tensiones, últimas etapas del canal, no tienen módulo propio y se exponen en el post-proceso —sonda puntual, vista tridimensional y tabla de resultados— y en la memoria de cálculo, que desarrolla las nueve etapas con sustitución numérica. A diferencia de las visualizaciones genéricas de los textos, los módulos operan directamente sobre el elemento seleccionado en el lienzo, conectando la abstracción matemática con la geometría concreta del problema: es la decisión de diseño que materializa el aporte pedagógico central del trabajo, hacer visible y manipulable lo que en la enseñanza tradicional permanece como una caja negra entre los datos de entrada y los resultados, y la matriz de trazabilidad la registra como encarnación de los principios de caja de cristal, contigüidad y señalización, segmentación, compromiso activo y constructivo y vínculo bidireccional entre la matriz y la geometría.

La verificación y validación del motor (\autoref{cap:resultados}) se abordó mediante una batería de tres niveles, con criterios de aceptación fijados a priori en los propios guiones. La verificación de código por MMS, ya mencionada, confirmó las tasas de convergencia teóricas. La contrastación con referencias externas se realizó con la viga de Timoshenko, contrastando las tensiones y la deflexión contra la solución elástica de Timoshenko y Goodier y contra un modelo equivalente en el software comercial SAP2000: frente a la solución analítica los errores fueron del $0{,}04\,\%$ en la tensión normal $\sigma_x$, del $0{,}26\,\%$ en la flecha y de hasta el $2{,}9\,\%$ en la tensión cortante; frente a SAP2000, del $0{,}21\,\%$ en $\sigma_x$ y de hasta el $0{,}56\,\%$ en desplazamientos; y la suma de las reacciones equilibra la carga total a precisión de máquina. Finalmente, la membrana de Cook reprodujo ese caso de referencia clásico y permitió ilustrar empíricamente el bloqueo por cortante (\emph{shear-locking}) del elemento Q4 frente a la convergencia rápida del Q9: para igual número de grados de libertad (GDL), el Q9 alcanzó un error de $-0{,}044\,\%$ donde el Q4 conservaba aún un $-0{,}594\,\%$. El marco metodológico de verificación y validación adoptado sigue las prácticas establecidas en la literatura de computación científica, con la salvedad, declarada en el \autoref{cap:marco_teorico}, de que al no disponerse de campaña experimental propia la contrastación se apoya en soluciones analíticas y en un modelo equivalente en un programa comercial; la validación en sentido estricto, contra mediciones físicas, queda fuera del alcance de este trabajo.

En cuanto a la memoria de cálculo y la interoperabilidad (\autoref{cap:diseno}; resultado en la \autoref{sec:cobertura} y en la \autoref{sec:consistencia-interna}), EduFEM genera una memoria de cálculo en PDF que reconstruye el procedimiento numérico paso a paso —de las funciones de forma a las tensiones nodales— con las mismas funciones del motor, generada sin error para Q4 y Q9 y reproducida en el \autoref{anexo:memoria}, y ofrece el intercambio de datos mediante importación de geometría DXF y modelos en formato CSV, verificado por ida y vuelta sin pérdida y por idempotencia de la importación. La memoria es, en los términos de la matriz de trazabilidad, la encarnación del efecto del ejemplo resuelto: integra las partes del procedimiento en un solo documento, rotula y segmenta sus pasos y los desarrolla sobre el problema del propio alumno. Esta capacidad cierra el ciclo de uso de la herramienta: el estudiante no solo experimenta con el modelo, sino que también documenta y reproduce sus cálculos.

Por último, en cuanto a la evaluación del potencial didáctico (\autoref{cap:resultados}, \autoref{sec:cobertura}), los tres instrumentos fijados a priori dieron el resultado siguiente. La tabla de especificaciones muestra que las once unidades de contenido del canal de cálculo tienen instrumento en los niveles recordar, comprender y aplicar; en el nivel analizar lo tienen ocho de once y en el nivel evaluar siete de once, y las celdas vacías se concentran en tres unidades —funciones de forma y mapeo, matriz de deformación-desplazamiento y fuerzas nodales equivalentes— en las que el software muestra el cálculo pero no permite variar un parámetro y observar su efecto. El mapa de conjeturas y su matriz de trazabilidad muestran que los diez principios de aprendizaje tienen encarnación verificable: ocho con veredicto cumple y dos con veredicto cumple parcialmente, el compromiso interactivo entre pares, ajeno a un software de escritorio individual, y el andamiaje de las simulaciones, que en EduFEM es fijo y no adaptativo. La evaluación heurística con LORI documentó evidencia en las ocho dimensiones aplicables: cuatro cumplidas —calidad del contenido, alineación con los objetivos, diseño de la presentación y reusabilidad—, tres parciales —retroalimentación y adaptación, usabilidad de la interacción y accesibilidad—, una no evaluable sin usuarios —motivación— y una no aplicable —cumplimiento de estándares de objetos de aprendizaje empaquetados—. El séptimo objetivo se cumple, además, con el diseño de la prueba de campo que las recomendaciones detallan.

En conjunto, los resultados confirman que EduFEM satisface los siete objetivos específicos en los términos en que fueron enunciados —el primero como fundamento y no como resultado medido; el cuarto con módulos hasta el ensamblaje y las dos últimas etapas expuestas en el post-proceso y la memoria; el séptimo con la evaluación del potencial y el diseño de la prueba de campo— y, con ello, el objetivo general. Confirman también, cláusula por cláusula, la parte de la hipótesis que este ciclo contrasta: la exposición del canal es relevante, porque cubre las once unidades de contenido en comprender y aplicar; es correcta, porque el motor reprodujo las tasas teóricas de convergencia con errores acotados frente a las referencias y evidenció el bloqueo por cortante del Q4 frente al Q9; es consistente, porque cada principio de aprendizaje tiene encarnación verificable y cada etapa del canal es observable en un módulo o en la memoria; y es practicable, porque ninguna dimensión aplicable del instrumento heurístico quedó sin evidencia. Que esa exposición apoye la comprensión de los fundamentos del método —la efectividad— es la conjetura teórica del mapa de conjeturas: fundamentada en el \autoref{cap:marco_teorico}, no contrastada en este ciclo y remitida a la prueba de campo del segundo.

\section*{Aportes del trabajo}

El principal aporte de este trabajo es una herramienta de software libre, funcional y verificada, que articula tres capas habitualmente disociadas en los recursos educativos sobre el MEF: un motor de cálculo riguroso, una interfaz de modelado interactiva y un conjunto de módulos que explican el método sobre el modelo real del estudiante.

Desde el punto de vista numérico, el motor implementa de manera completa y verificada la formulación isoparamétrica Q4 y Q9 para los dos estados planos de la elasticidad, con un esquema de indexación de GDL robusto frente a identificadores de nodo no contiguos —que desacopla la identidad del nodo de su índice de ensamblaje— y con tratamiento de cargas volumétricas, desplazamientos prescritos no homogéneos y normas de error $L^2$ y $H^1$ integradas con cuadratura de orden adecuado. La validación acota el error en las magnitudes primarias por debajo del $0{,}3\,\%$ frente a la solución analítica y del $0{,}6\,\%$ frente al modelo equivalente en SAP2000, para los problemas de referencia estudiados.

Desde el punto de vista pedagógico, el aporte distintivo es la decisión de diseño de que cada módulo educativo opere sobre el elemento que el estudiante selecciona en el lienzo compartido, en lugar de sobre un ejemplo prefabricado. Esto convierte la exploración del Jacobiano, de la matriz $\bm{B}$ o de la cuadratura de Gauss en una experiencia ligada al problema concreto, y permite que fenómenos sutiles como el bloqueo por cortante o los modos espurios (\emph{hourglass}) se observen en el modelo propio. La memoria de cálculo en PDF eleva esta capa a un segundo aporte pedagógico: genera, sobre el modelo concreto del usuario, un documento que reconstruye el procedimiento completo con sustitución numérica paso a paso —de las funciones de forma a las tensiones nodales—. Frente a las visualizaciones interactivas, que muestran \emph{cómo} se calcula, la memoria deja un registro escrito que el alumno puede seguir y reproducir por sus propios medios, cerrando el lazo entre la teoría, la exploración y el cálculo manual.

Desde el punto de vista metodológico, el trabajo aporta el mapa de conjeturas y los tres instrumentos documentales con que evaluó el potencial didáctico —la tabla de especificaciones sobre las once unidades del canal de cálculo, la matriz de trazabilidad sobre diez principios de aprendizaje y la evaluación heurística con LORI—, reutilizables para evaluar otros recursos de enseñanza del MEF y para definir qué debe observar la prueba de campo del segundo ciclo: sus celdas vacías y sus veredictos parciales son, a la vez, las limitaciones del artefacto y las preguntas de esa prueba.

\section*{Limitaciones}

Las fronteras del trabajo se fijaron a priori en la Introducción (\emph{Alcance y limitaciones}): elasticidad lineal estática en dos dimensiones —sin sistemas estructurales discretos, placas, cáscaras ni problemas tridimensionales—, elementos cuadriláteros Q4 y Q9, material elástico lineal e isótropo, malla construida manualmente o por generación estructurada, y evaluación de la dimensión educativa limitada al potencial didáctico. A ellas se suman las que el desarrollo y los resultados pusieron de manifiesto. En primer lugar, el potencial didáctico no es efectividad: los tres instrumentos establecen que el artefacto está en condiciones de producir el efecto buscado, no que lo produzca, y la evaluación heurística es una autoevaluación documentada del autor, no un juicio independiente. En segundo lugar, la cobertura del canal por módulos interactivos llega hasta el ensamblaje: la solución del sistema y la recuperación de tensiones se exponen en el post-proceso y en la memoria, y no en un módulo propio; y tres unidades de contenido —funciones de forma y mapeo, matriz de deformación-desplazamiento y fuerzas nodales equivalentes— no ofrecen la manipulación de un parámetro que sostendría los niveles cognitivos de analizar y evaluar. En tercer lugar, el elemento Q4 exhibe bloqueo por cortante bajo flexión dominante, fenómeno que el software ilustra deliberadamente con fines didácticos pero que no mitiga mediante técnicas correctivas, y en deformación plana la matriz constitutiva degenera hacia la incompresibilidad cuando $\nu \to 0{,}5$, régimen que el módulo M4 señala pero que el software no resuelve. En cuarto lugar, el solucionador emplea una factorización directa sobre la matriz de rigidez almacenada en formato disperso, esquema eficiente para los tamaños de modelo propios de un contexto educativo ---e incluso para mallas de decenas de miles de grados de libertad, según cuantifica el \autoref{cap:resultados}--- pero cuyo coste crece de forma desfavorable en el extremo de mallas masivas. En quinto lugar, la validación descansa en tres casos de referencia y en un único caso con material de ingeniería civil; una batería más amplia reforzaría la confianza ante geometrías y cargas más diversas, y la carga superficial linealmente variable no interviene en ningún caso de referencia. Finalmente, la paleta de color arcoíris de los contornos, elegida por continuidad con la herramienta profesional, limita la accesibilidad para la visión cromática deficiente, como la evaluación heurística registra.

\section*{Recomendaciones y trabajo futuro}

A partir de las limitaciones anteriores se plantean tres líneas de continuación, una por capítulo, cada una con sus ítems concretos.

\subsection*{Sobre la formulación (\autoref{cap:marco_teorico})}

\textbf{Tratamiento del bloqueo.} Para abordar el bloqueo (\autoref{sec:locking}) —por cortante en flexión y volumétrico cuando $\nu \to 0{,}5$ en deformación plana— se recomienda incorporar técnicas correctivas como la integración reducida selectiva (SRI) o la formulación $\bar{\bm{B}}$ (B-bar) \autocite{hughes2000fem,bathe2014fem}; conviene recordar que el Q9 con integración completa tampoco está libre del bloqueo volumétrico, de modo que estas técnicas beneficiarían a ambos elementos. Más allá de su valor numérico, estas técnicas tienen un alto valor pedagógico: permitirían un módulo educativo dedicado que contraste, sobre un mismo modelo, el comportamiento bloqueado y el corregido, profundizando en la comprensión del fenómeno que la herramienta ya ilustra con el caso de referencia de Cook.

\textbf{Más tipos de elemento.} La biblioteca de elementos puede ampliarse con elementos triangulares (lineales y cuadráticos) y con cuadriláteros de transición, ofreciendo así una variedad que enriquezca el estudio comparativo de la convergencia y del comportamiento bajo distorsión de malla \autocite{onate2009structural,reddy2006introduction}.

\subsection*{Sobre el software (\autoref{cap:diseno})}

\textbf{Manipulación en las unidades sin nivel de análisis.} Las celdas vacías de la tabla de especificaciones señalan la extensión de mayor rendimiento didáctico por unidad de esfuerzo: permitir que el alumno desplace un punto en coordenadas naturales y vea su imagen física en M1, distorsione el elemento y observe el cambio de la matriz $\bm{B}$ en M3, y modifique la intensidad o la distribución de una carga superficial y vea las fuerzas nodales en M6. Cada una de esas manipulaciones convertiría una unidad que hoy solo se comprende y se aplica en una que también se analiza, en los términos de la taxonomía revisada de Bloom \autocite{anderson2001taxonomy}.

\textbf{Solucionador para problemas de gran porte.} El motor ya ensambla la matriz de rigidez global vectorizada por lotes, la almacena en formato disperso comprimido (CSR) y resuelve el sistema reducido con una factorización LU dispersa directa con ordenamiento de mínimo grado, cuyo efecto cuantifica el \autoref{cap:resultados}. Las extensiones de mayor impacto sobre la escalabilidad consisten en incorporar una factorización de Cholesky —aplicable porque $\bm{K}_{ff}$ es simétrica y definida positiva (\autoref{sec:resolucion-tensiones})— y en ofrecer un solucionador iterativo precondicionado para problemas masivos. Estas mejoras preservarían la interfaz pública del solucionador y la versión legible del ensamblaje reservada a los módulos educativos, de modo que no obligarían a reescribir las capas de post-proceso ni dichos módulos.

\textbf{Extensión a tres dimensiones.} A más largo plazo, la extensión del motor a problemas tridimensionales con elementos hexaédricos (H8, H20) y tetraédricos constituiría un salto natural, aprovechando que gran parte de la formulación isoparamétrica y de la infraestructura de cuadratura es generalizable a tres dimensiones.

\textbf{Mallado automático y accesibilidad.} Un generador de mallas sobre un contorno dibujado en el lienzo, con las métricas de calidad del módulo M0 como control, reduciría el trabajo manual que hoy exige un modelo de más de unas decenas de elementos y ampliaría los casos que el estudiante puede plantear por sí mismo. Una paleta perceptualmente uniforme, ofrecida como alternativa a la paleta arcoíris, atendería la dimensión de accesibilidad que la evaluación heurística dejó en veredicto parcial.

\subsection*{Sobre la evidencia (\autoref{cap:resultados})}

\textbf{Ampliación de la batería de validación.} La evidencia del \autoref{cap:resultados} descansa en tres casos de referencia y en un único caso con material de ingeniería civil, la viga de hormigón. Se recomienda incorporar un caso con carga superficial linealmente variable, que ningún caso de referencia ejercita, y al menos un caso civil en deformación plana con solución de referencia publicada —una presa de gravedad, un muro de contención o un túnel somero— con peso propio, de modo que la validación externa cubra el estado plano que la práctica civil emplea con más frecuencia.

\textbf{Segundo ciclo: prueba de campo de la efectividad.} La línea de continuación más importante para consolidar la finalidad del proyecto es la prueba de campo que contraste la conjetura teórica del mapa de conjeturas: que la exposición del canal apoye la comprensión de los fundamentos del método. El primer ciclo deja ese diseño especificado en sus elementos esenciales. La población es la de la \autoref{sec:poblacion}, los estudiantes de la Carrera de Ingeniería Civil de la Universidad Autónoma Tomás Frías que cursan la asignatura en que se introduce el MEF, y la muestra, dirigida, un curso completo de esa asignatura. El diseño es de preprueba y posprueba con un solo grupo, o con dos cohortes si el calendario académico lo permite, en la línea de la comparación entre cohortes con que Bishay midió el efecto de sus herramientas \autocite[pp.~4-5, 13]{bishay2020teaching}. El tratamiento reproduce las estructuras de tarea del mapa de conjeturas: construir el ejemplo canónico, validar su salud, resolverlo, recorrer los módulos M1 a M7 sobre un elemento propio, generar la memoria y reproducir a mano una etapa, y convertir el modelo de Q4 a Q9 y refinar la malla sobre la membrana de Cook. El instrumento es una prueba de conocimientos sobre las once unidades de contenido, con ítems por nivel cognitivo tomados de la tabla de especificaciones, complementada por un cuestionario de percepción para las dimensiones de LORI que este ciclo no pudo evaluar —motivación y usabilidad de la interacción—, en la línea de la evaluación cualitativa que Lee aplicó en sus cursos \autocite[p.~168]{lee2015interactive}. El análisis compara la preprueba y la posprueba con estadística no paramétrica para muestras pareadas, adecuada al tamaño de un curso, y los procesos mediadores del mapa se observan en los artefactos que cada alumno produce: sus archivos de proyecto y las memorias generadas y reproducidas. Los resultados de esa prueba permitirían afinar el diseño de los módulos educativos sobre evidencia y aportarían el sustento empírico a la utilidad didáctica que el presente ciclo fundamenta por diseño y por la literatura.

En conjunto, estas líneas trazan un camino de evolución que preserva la identidad educativa de EduFEM al tiempo que amplía su alcance numérico, su biblioteca de elementos y, sobre todo, su fundamentación pedagógica empírica.
